\chapter{Introduction}
\thispagestyle{empty}

\section{Background}
Neurodevelopmental disorders are conditions that begin in the developmental period and affect
personal, social, academic or occupational functioning~\cite{apa2013}. Two of them are studied
in this project. Attention-deficit/hyperactivity disorder (ADHD) is characterised by
inattention, hyperactivity and impulsivity; a pooled estimate of its worldwide prevalence in
children and adolescents is about 5\%~\cite{polanczyk2007}. Intellectual and developmental
disorder (IDD), also called intellectual disability, combines significantly below-average
intellectual functioning with limitations in adaptive behaviour; its prevalence is about
1\%~\cite{maulik2011}.

Both disorders are diagnosed from behaviour: clinical interviews, rating scales completed by
parents and teachers, and standardised tests, interpreted by a specialist according to criteria
such as DSM-5~\cite{apa2013}. The process is slow, depends on the availability of specialists
and involves subjective judgement. An objective physiological measurement that supports the
diagnosis would therefore be valuable.

Electroencephalography (EEG) records the electrical activity of the brain from electrodes on
the scalp. It is non-invasive, inexpensive compared with other brain-imaging methods, portable
and has millisecond time resolution, and consumer-grade wireless headsets make recordings
possible outside the hospital. These properties make EEG the most practical candidate for a
physiological aid to the assessment of neurodevelopmental disorders.

\section{Problem statement}
The problem addressed is the \textbf{detection of multiple neurodevelopmental disorders from
EEG}: given an EEG recording, decide whether it comes from a person with the disorder or from a
control, for IDD and for ADHD, with one common processing and classification framework.

Two requirements follow from the intended use. The framework should handle more than one
disorder with the same processing, rather than a method tuned separately for each condition. And
it must work on a person whose EEG it has never seen, because that is the only situation in
which a detection system is used. The second requirement decides how every result in this
report is obtained.

\section{Motivation for the proposed pipeline}
Published methods report accuracies of 95--100\% on the two public datasets used here
(Chapter~2), but most of them divide the short windows of each person between the training and
the test set, or do not report keeping people apart, so that the classifier may be tested on
people it has already seen. When the
method with the highest reported accuracy was re-evaluated on unseen people in this project, its
accuracy fell to little above chance (Chapter~4). A broad comparison of alternatives under a
strict protocol (Chapter~5) then showed that representing each window by the
\emph{covariance matrix of its channels}, analysed with Riemannian geometry, carried the most
information about ADHD. The proposed pipeline (Chapter~6) builds on this finding: it keeps the
adaptive multichannel decomposition (SMVMD) of the starting method to split the signal into
frequency-specific common modes, describes each mode by its channel covariance, and classifies
the Riemannian tangent-space features with gradient-boosted trees.

\section{Objectives}
\begin{enumerate}
\item to develop an EEG classification pipeline that applies unchanged to IDD and ADHD and is
      light enough to run on an ordinary computer without a GPU;
\item to evaluate it on people not seen during training, with a protocol that excludes every
      form of information leakage;
\item to analyse each component of the pipeline (decomposition, common modes, Riemannian
      representation, classifier) by ablation, and to compare the result with published work;
\item to establish what the pipeline relies on, which modes and electrodes carry the
      information, and which properties of the data other than the disorder could explain the
      results.
\end{enumerate}

\section{Contributions}
\begin{itemize}
\item \textbf{A pipeline combining SMVMD common modes, Riemannian geometry and gradient
      boosting.} The window-specific SMVMD modes are mapped to four common modes by centre
      frequency, so that a fixed set of covariance matrices describes every window. On unseen
      ADHD children it reaches \R{MainDStwoBal}\% balanced accuracy and \R{MainDStwoSubj}
      children correct; on IDD all participants are classified correctly.
\item \textbf{A strict evaluation} with subject-wise cross-validation, nested tuning on
      training subjects only, an automated leakage test and majority voting per person.
\item \textbf{Ablations} of the classifier (LR, SVM, XGBoost, LightGBM), of the individual
      modes, and of the decomposition itself against fixed band-pass filters and the undecomposed
      signal.
\item \textbf{A quantified comparison with the literature}, showing how window-level splitting
      inflates reported accuracy, and that the proposed pipeline is comparable with the best
      published subject-wise ADHD result without a deep network or a GPU.
\end{itemize}

\section{Organisation of the report}
Chapter~2 reviews EEG-based detection of ADHD and IDD and the limitations of published
evaluations. Chapter~3 describes the datasets and their preprocessing, and Chapter~4 the
evaluation protocol. Chapter~5 summarises the preliminary study that led to the proposed
pipeline. Chapter~6 describes the pipeline in detail, and Chapter~7 presents its results,
ablations and analyses. Chapter~8 gives the conclusions, limitations and future work.
